\section{The study design}\label{ch:design}

The experiments have to show three things: that the losses of \cref{ch:diagnosis} occur, that FedECS-LT recovers
them, and which part of FedECS-LT does so. Every choice below serves one of these, or protects the comparison against
a way of being unfair.

\subsection{The research questions}\label{sec:s-rq}

\begin{itemize}
\item \textbf{RQ1, diagnosis.} Do the three losses occur in GeFL-F's own runs? Answered by the class rows, the readout
gain and the logit spread (\cref{ch:diagnosis}).
\item \textbf{RQ2, effect.} Does FedECS-LT, ending with its readout, raise balanced accuracy over GeFL-F and the other
model-heterogeneous methods, each as it is without the readout? Answered by the main comparison (\cref{sec:s-main}).
\item \textbf{RQ3, body or last layer.} How much of the difference remains when every method of the GeFL-F pipeline
gets the same readout, so that only the heads' bodies differ? Answered by the readout comparison
(\cref{sec:s-same-readout}).
\item \textbf{RQ4, scope.} Does the effect hold with natural class imbalance and in the protocol with one shared
model? Answered in \cref{sec:s-natural,sec:s-shared}.
\item \textbf{RQ5, components.} Which changes carry the effect? Answered by the ablations, whose runs are pending
(\cref{sec:s-pending}).
\end{itemize}
RQ3 is the question that decides how the paper's claim must be worded, as \cref{ch:results} shows.

\subsection{Datasets and splits}

\paragraph{The long-tailed benchmarks} CIFAR-10-LT and CIFAR-100-LT are built by the paper with $\rho=100$. GeFL-F
uses half of CIFAR-10's training set, so the profile is applied to that half (and to 22\% of CIFAR-100, which gives the
usual CIFAR-100-LT shape). As \cref{ch:problem} showed, the largest classes are capped, so the realised ratios are 48
and 96 (\cref{tab:s-datasets}). The class index orders classes from the most to the least frequent, so the tail
classes are the same in every run. MNIST, Fashion-MNIST and SVHN belong to the same benchmark set; their runs had not
returned when this text was written.

\paragraph{The naturally imbalanced datasets} Six datasets keep their own class sizes, with training sets capped at
20,000 images: DermaMNIST, OrganAMNIST, BloodMNIST and PathMNIST (medical images), EMNIST ByClass (handwritten
characters, 62 classes) and FER-2013 (facial expressions). Their realised ratios range from 2 to 75. On OrganAMNIST,
BloodMNIST and PathMNIST the ratio is 5 or less, so there the results speak to label skew more than to a long tail.
These six datasets were never used while the method was being developed, which makes them the cleanest evidence in
the paper.

\paragraph{Splits} Each class is divided among $K=10$ clients by a Dirichlet draw with $\alpha=0.5$. The draw is
repeated, at most 100 times, until every client holds at least ten images. A held-out block of training images that
no client holds, 200 per class on CIFAR-10, forms the server's validation set. It has two uses only: it trains the
referee of the fidelity measurement, and it sets the public bounds of secure aggregation. It is never used to select a
model or a setting. Every reported number is on the official test set. The paper checked that every client's class
counts are identical across methods for every dataset and seed.

\paragraph{Seeds} Every configuration runs with seeds 3, 4 and 5. A seed fixes the split and every method's
initialisation, so two methods with the same seed see exactly the same clients: their results are \emph{paired}.

\subsection{The two protocols}

\paragraph{Different client models} The main setting is GeFL-F's: 10 clients, one per architecture, with GeFL-F's
rounds ($T_{\mathrm I},T_{\mathrm{II}},T_{\mathrm{III}}=50,200,100$ on CIFAR). All federated methods of this setting
share the same Stage-I extractor for a given dataset and seed, so they differ only in what they do after Stage I.

\paragraph{One shared model} The second protocol is the one used by federated long-tail research, taken from the
FedLTLib benchmark, whose settings come from CReFF: CIFAR-10-LT and CIFAR-100-LT at $\rho=100$, 20 clients, 8 of them
(40\%) in each round, Dirichlet $\alpha=0.5$, one ResNet-8 shared by all, 200 rounds of 10 local epochs (batch 32, SGD
with learning rate 0.1, momentum 0.9, decay 0.99 per round). The baselines run their authors' LongTailLib code
unchanged on the paper's splits, with ResNet-8 for every method.

FedECS-LT is fitted into this protocol by splitting the ResNet-8 after its second stage into a frozen extractor and a
head. The extractor trains for 100 rounds, the generator for 100, and the head for 100 with 10 real epochs per round;
so FedECS-LT uses the baselines' 200 training rounds plus 100 generator rounds. Two differences matter later: its head
trains on frozen features of unaugmented images, while the baselines train the whole network end to end with
augmentation; and with a single architecture the consensus soft label is just the shared model's own prediction.

\subsection{The methods compared}\label{sec:s-methods}

Each group of methods answers a specific question.
\begin{itemize}
\item \textbf{GeFL-F}, the method FedECS-LT extends: the difference to it is the effect of the changes.
\item \textbf{FedDF}, ensemble distillation on the server: after each round the server distils every architecture
toward the ten architectures' mean prediction, for 20 steps. FedDF originally distils on public data; here it distils
on GeFL-F's generated features, so that no method uses data that no client holds. It is the strongest generator-based
competitor.
\item \textbf{Per-client long-tail losses}: Balanced Softmax, FedLC and FedRS, applied to the heads with \emph{no
generator}. They test whether correcting the loss alone repairs the long tail. By the reparametrisation lemma they
should behave alike, and \cref{ch:diagnosis} showed that they do once the readout is added.
\item \textbf{Model-heterogeneous methods from the HtFLlib library}: prototypes (FedProto, FedTGP), a generator of
embeddings (FedGen), a shared global head (FedGH), shared last layers (LG-FedAvg) and shared sub-networks (FedMRL). They
represent the other ways of sharing knowledge between different models. They were designed for label skew without a
global long tail and are ported to the frozen extractor with HtFLlib's settings.
\item \textbf{Local}: each client trains its extractor and head alone. It is the reference without federation.
\item \textbf{One-shared-model methods}: FedAvg, FedProx, MOON, FedLC, CCVR, CReFF, FedETF, Fed-GraB, FedLoGe, FedYoYo
and RUCR, in the second protocol only.
\end{itemize}
Two further groups, GeFL's own grouped baselines and a reference trained on all data pooled, are planned for GeFL's
50-client setting and had not returned when this text was written.

\paragraph{Fairness of the comparison} The paper's fairness table records what each method spends and exchanges.
FedECS-LT trains 10 synthetic epochs per round against 1 for GeFL-F and FedDF and none for the methods without a
generator; it exchanges consensus soft labels; and it ends with the readout. GeFL-F, FedDF and the three losses are
also reported with the readout added. The HtFLlib methods are not, so FedECS-LT's margin over them includes the
readout's gain. No baseline was tuned on the paper's data.

\subsection{Metrics}

The primary metric is balanced accuracy averaged over the ten client models (\cref{ch:problem}). The others each answer
a narrower question:
\begin{itemize}
\item \textbf{tail recall}, the mean recall of the 40\% least frequent classes (classes 6--9 of CIFAR-10): are the rare
classes recovered?
\item \textbf{worst-class recall}, the lowest class recall after averaging over the models: is any class abandoned?
\item \textbf{classes with zero recall for every model}: how many classes does the federation fail completely?
\item \textbf{accuracy on each client's own class mix}: the test accuracy re-weighted to the class frequencies of the
model's own client. It measures how well a model serves its own client, which balanced accuracy deliberately ignores.
\item \textbf{expected calibration error} (15 bins): do the predicted probabilities match the observed frequencies?
\item \textbf{the weakest architecture's balanced accuracy}: does the mean over ten models hide a poor one?
\item \textbf{GeFL's measure}, each model's best test accuracy over the evaluated rounds, given in the appendix for
comparison with GeFL's paper. The main tables report the model after the last round, which involves no selection on
the test set.
\item \textbf{cost}: communication per client, wall-clock time (on a shared T4 GPU, so only indicative) and
floating-point operations.
\end{itemize}

\subsection{Statistics}\label{sec:s-stats-design}

\paragraph{Paired differences} A difference between FedECS-LT and a method is the mean of the per-seed differences
over the seeds both have, with a 95\% $t$ interval and a paired $t$-test. Pairing removes the variation that comes from
the split, because both methods see the same split on each seed.

A worked example with the run logs: with the readout, FedECS-LT minus Balanced Softmax on CIFAR-10-LT is $0.65$,
$1.30$ and $2.34$ points on seeds 3, 4 and 5. The mean is $1.43$ and the standard deviation $0.86$. The $t$ statistic
is $1.43/(0.86/\sqrt3)=2.89$ with 2 degrees of freedom, giving an uncorrected two-sided $p=0.10$. The 95\% interval is
$1.43\pm4.30\cdot0.86/\sqrt3$, which is \R{main/cifar10/er_ci/Balanced Softmax}: it contains zero.

\paragraph{Holm's correction} A table compares FedECS-LT with many methods at once, and with many tests some will look
significant by chance. Holm's procedure controls this. Sort the $m$ $p$-values from smallest to largest, multiply the
smallest by $m$, the next by $m-1$, and so on, and keep each adjusted value at least as large as the one before. A
difference is called significant only if its adjusted $p$ is below $0.05$. The family is every other method of the
same table and setting. In the example, after correction over the five methods of the pipeline, Balanced Softmax's
$p$ becomes \R{main/cifar10/er_p/Balanced Softmax}.

\paragraph{Why three seeds are weak} Three paired differences give little power. A sign-flip permutation test, which
assumes nothing about the distribution, cannot go below $p=0.25$ with three pairs: of the $2^3=8$ ways to flip the
signs, the observed pattern and its mirror are always 2 of the 8. The $t$-test reaches smaller values only because it
assumes normally distributed differences. A result near $p=0.05$ can also depend on which methods are in the Holm
family.

\paragraph{The noise floor} GPU computation is not bit-reproducible, so two runs with the same seed differ. The paper
measured this by running the BloodMNIST notebook twice with the same run list (\R{repeat/bal/n} runs). The standard
deviation of the difference between the two runs of one configuration is \R{repeat/bal/sd} points as trained and
\R{repeat/er/sd} with the readout. For a mean of three paired differences that gives about half of these values, and
the paper treats differences smaller than twice that, \R{repeat/bal/floor} points as trained and \R{repeat/er/floor}
with the readout, as unresolved. Exceeding the floor is necessary, not sufficient: a difference above it counts as
resolved only if its Holm-corrected test is also significant.

\subsection{What the development saw}\label{sec:s-development}

The settings of FedECS-LT were fixed during its development, before the reported runs. The development used seeds 0 to
5, including seeds 3 to 5 with the same client splits as the reported runs, and every run was evaluated on the same
official test sets. It compared about 36 configurations of the method and its variants on CIFAR-10-LT and 17 on
CIFAR-100-LT, plus runs on MNIST, Fashion-MNIST and SVHN. It never used the six naturally imbalanced datasets. The
baselines were not tuned. So on the CIFAR benchmarks the reported numbers may favour FedECS-LT through this selection,
which is why the natural datasets matter, and why a run on new seeds 6--8 was added (pending).

\subsection{Threats to validity}

The paper lists these threats in its design section. Each one changes how a result may be read.
\begin{itemize}
\item \textbf{Compute.} FedECS-LT trains longer and communicates more. Its gain is therefore not shown to come from the
statistics rather than from the extra computation until the matched-epoch runs return.
\item \textbf{Missing runs.} Some settings have fewer than three seeds (EMNIST, and CIFAR-100-LT with one shared model),
and those values carry no interval. With one shared model some runs of FedECS-LT failed where the baselines finished,
so the missing runs are not random.
\item \textbf{Ports.} The HtFLlib methods run on a frozen extractor with settings meant for end-to-end training, which
may not suit them.
\item \textbf{The mean over models.} Balanced accuracy averages ten models of different sizes, so a large gain on small
heads and a small one on large heads give the same number; the weakest architecture is reported for this reason.
\item \textbf{The readout comparison.} The readout is the same linear classifier on every body, so accuracy with it
measures the bodies only through that classifier.
\item \textbf{The validation set.} The server holds a small validation set for the diagnostics and the bounds; a
deployment without such data would need another source for the bounds.
\item \textbf{The long tail is synthetic} on the benchmarks, with the same tail classes in every run.
\end{itemize}

What do the results show?
